\section{Zero exclusion with explicit positive margins}
\label{sec:elementary-zero-exclusion}

The distinction between a negative statement and a positive numerical
estimate matters here.  A proof of $\zeta(s)\ne0$ obtained by contradiction
does not, by itself, supply a primitive recursive precision at which a
numerical computation will certify nonvanishing.  We shall prove positive
rational lower bounds.  They are much weaker than the best analytic
bounds, but their role in the arithmetic construction is transparent.

We use the continuation formula
\begin{equation}
 I(s)=\int_1^\infty\{x\}x^{-s-1}\,dx,
 \qquad \zeta(s)=\frac{s}{s-1}-sI(s)
 \quad(\Re s>0,\ s\ne1).
 \label{eq:elem-euler-integral}
\end{equation}
Here and below an integral is shorthand for its finite quadrature
approximants and their stated tail estimates.  The braces in the
integrand cause no problem: split the integral at integer endpoints,
where its pieces are the smooth functions $(x-n)x^{-s-1}$.
For $\Re s\ge a>0$,
\begin{equation}
 |I(s)|\le a^{-1},\qquad |I'(s)|\le a^{-2},\qquad
 \left|\int_X^\infty\{x\}x^{-s-1}\,dx\right|
 \le \frac{X^{-a}}a .
 \label{eq:elem-I-bounds}
\end{equation}
The derivative estimate follows by inserting $\log x$ and using
$\int_1^\infty (\log x)x^{-a-1}dx=a^{-2}$; finite integration by parts
proves this last identity with its exponentially decreasing remainder.
Differentiating the finite approximants and their remainder bounds gives
\begin{equation}
 \zeta'(s)=-\frac1{(s-1)^2}-I(s)-sI'(s).
 \label{eq:elem-zeta-derivative}
\end{equation}

\subsection{The Euler region}

\begin{theorem}[Positive Euler-region nonvanishing]
\label{thm:elem-euler}
For $s=\sigma+it$ with $\sigma>1$,
\begin{equation}
 |\zeta(s)|\ge\frac{\sigma-1}{\sigma}.
 \label{eq:elem-euler-apartness}
\end{equation}
In particular, a supplied rational gap $\sigma\ge1+1/q$, $q\ge1$,
gives the rational lower bound $1/(q+1)$.
\end{theorem}

\begin{proof}
Define $D(s)=\sum_{n\ge1}\mu(n)n^{-s}$.  Trial division computes
$\mu(n)$ by primitive recursion.  Expanding the product over the distinct
prime divisors of $n$ gives the finite identity
$\sum_{d\mid n}\mu(d)=[n=1]$.  Consequently
\[
 \zeta(s)D(s)=1.
\]
For clarity, the passage from the finite convolution to this equality
has an explicit bound.  A product tail with $ab>N^2$ is contained in the
union $a>N$ or $b>N$.  Its absolute value is at most
\[
 2\zeta(\sigma)\frac{N^{1-\sigma}}{\sigma-1}.
\]
Thus the divisor calculation and this estimate prove the asserted
identity at every requested precision.  Since
\[
 |D(s)|\le\zeta(\sigma)
 \le1+\int_1^\infty x^{-\sigma}\,dx
 =\frac{\sigma}{\sigma-1},
\]
the result follows by multiplying the two norms.  No test for equality
of a complex number to zero is involved.
\end{proof}

\subsection{The finite 3--4--1 inequality}

\begin{lemma}[Euler positivity]
\label{lem:elem-341}
For $\sigma>1$ and real $t$,
\begin{equation}
 \zeta(\sigma)^3\,|\zeta(\sigma+it)|^4
                 |\zeta(\sigma+2it)|\ge1.
 \label{eq:elem-341}
\end{equation}
\end{lemma}

\begin{proof}
Put
\[
 L(s)=\sum_p\sum_{k\ge1}\frac{p^{-ks}}k.
\]
The sum restricted by $p^k\le N$ has error at most
$N^{1-\sigma}/(\sigma-1)$: each prime power is a distinct integer, and
its coefficient $1/k$ is at most one.  The power series for
$-\log(1-z)$, applied to finite Euler products, gives
$\exp L(s)=\zeta(s)$.  This equality also has an elementary tail proof:
expand each finite geometric product, and observe that the omitted
integers have a prime divisor beyond the included prime bound and hence
exceed that bound.

For every finite prime-power sum, use
\[
 3+4\cos u+\cos(2u)=2(1+\cos u)^2\ge0
\]
with $u=kt\log p$, multiplied by $p^{-k\sigma}/k$.  Summing and then
using the explicit tails gives
\[
 3L(\sigma)+4\Re L(\sigma+it)+\Re L(\sigma+2it)\ge0.
\]
Exponentiation proves \eqref{eq:elem-341}.  Equivalently, at precision
$2^{-m}$ the finite statement has lower bound $1-2^{-m}$ after the
error budget for exponentiation and multiplication is included.  For a
supplied gap $\sigma-1\ge1/q$, the choice
$N=2^{q(m+q)}$ already bounds the logarithm tail by $2^{-m}$, because
$q2^{-q}\le1$.  Enlarging $m$ by the explicitly computed multiplication
budget completes the finite version.
\end{proof}

\subsection{A uniform zero-free neighborhood of the one-line}

Let $F(s)=(s-1)\zeta(s)$, with the removable value $F(1)=1$.

\begin{theorem}[A rational uniform clearance bound]
\label{thm:elem-F-clearance}
For every integer $H\ge1$, put
\[
 \varepsilon_H=2^{-40}(H+1)^{-5},\qquad
 b_H=2^{-36}(H+1)^{-4}.
\]
Then
\begin{equation}
 |F(\sigma+it)|\ge b_H
 \quad\text{if}\quad |t|\le H,\quad
 |\sigma-1|\le\varepsilon_H.
 \label{eq:elem-F-clearance}
\end{equation}
In particular $\zeta$ has no zero on $\Re s=1$, and has no zero in
the displayed neighborhood of that line; its pole is represented by
the nonzero removable value of $F$.
\end{theorem}

\begin{proof}
First suppose $1/8\le|t|\le H$.  Set
\[
 M=128(H+1),\qquad B=16(H+1),\qquad
 h=(128BM^4)^{-1},\qquad\delta=Mh.
\]
For $1/2\le\sigma\le2$,
\eqref{eq:elem-I-bounds}--\eqref{eq:elem-zeta-derivative} give
\[
 |\zeta'(\sigma+it)|
 \le64+2+4(H+2)=4H+74\le M.
\]
For $1\le\sigma\le2$, the same integral formula gives
\[
 |\zeta(\sigma+2it)|
 \le(2+2H)(4+1)=10(H+1)\le B.
\]
Indeed $|\sigma+2it-1|\ge1/4$ and $|I(\sigma+2it)|\le1$.
Also $0<h\le1$ and
$\zeta(1+h)\le1+1/h\le2/h$.  Apply
Lemma~\ref{lem:elem-341} at $\sigma=1+h$ to obtain
\[
 |\zeta(1+h+it)|^4
 \ge\frac{h^3}{8B}=16M^4h^4=16\delta^4.
\]
The fourth power is increasing on nonnegative quantities, so
$|\zeta(1+h+it)|\ge2\delta$.  The derivative bound along the
horizontal segment now gives
\[
 |\zeta(1+it)|\ge2\delta-Mh=\delta.
\]
On the larger segment $|\sigma-1|\le h/2$ it gives
\begin{equation}
 |\zeta(\sigma+it)|\ge\delta-Mh/2=\delta/2.
 \label{eq:elem-zeta-near-one}
\end{equation}
The constants simplify to
\[
 \delta=2^{-32}(H+1)^{-4},\qquad h/2=\varepsilon_H.
\]

We next handle the pole and small ordinates directly.  From
\eqref{eq:elem-euler-integral},
\[
 F(s)=s-s(s-1)I(s).
\]
If $|s-1|\le r<1$, then
\[
 |F(s)-1|\le r+\frac{r(1+r)}{1-r}
 =\frac{2r}{1-r}.
\]
In particular
\begin{equation}
 |F(s)|\ge1/3\qquad (|s-1|\le1/4).
 \label{eq:elem-pole-disk}
\end{equation}
For $|t|\le3/16$ and $|\sigma-1|\le\varepsilon_H$, we have
$|s-1|\le1/4$, since $\varepsilon_H\le1/16$.  For
$|t|\ge1/8$, \eqref{eq:elem-zeta-near-one} instead gives
\[
 |F(s)|\ge\tfrac18\,\tfrac\delta2
 =\delta/16=b_H.
\]
These two alternatives overlap.  They can be selected constructively:
approximate $|t|$ with error at most $1/64$, and compare the rational
answer with $5/32$.  The smaller alternative implies $|t|<3/16$;
the larger implies $|t|>1/8$.  This is a bounded rational comparison,
not an invocation of trichotomy for arbitrary reals.  Both alternatives
give \eqref{eq:elem-F-clearance}.
\end{proof}

This theorem supplies an elementary polynomial-width zero-free region.
It does not assert the stronger classical width $c/\log(|t|+2)$.
The numerical constants were chosen to expose the arithmetic witnesses,
not to optimize the region.  On the simpler band $1\le|t|\le H$, the
same proof with $M=16(H+1)$ and $B=4(H+1)$ gives the better bounds
\[
 |\zeta(1+it)|\ge2^{-21}(H+1)^{-4},\qquad
 |\zeta(\sigma+it)|\ge2^{-22}(H+1)^{-4}
\]
whenever $|\sigma-1|\le2^{-26}(H+1)^{-5}$.

\section{Trivial zeros, the critical strip, and symmetry}
\label{sec:elementary-location}

We use the entire completion in its explicitly convergent form
\begin{align}
 \psi(u)&=\sum_{n\ge1}e^{-\pi n^2u},\notag\\
 J(s)&=\int_1^\infty\psi(u)
       \bigl(u^{s/2-1}+u^{-(s+1)/2}\bigr)\,du,\notag\\
 \xi(s)&=\frac12+\frac{s(s-1)}2J(s).
 \label{eq:elem-xi-integral}
\end{align}
The effective Gaussian identity and its Mellin calculation, established
in the analytic setup, identify this with
$\tfrac12s(s-1)\pi^{-s/2}\Gamma(s/2)\zeta(s)$ in $\Re s>1$.
The displayed integral itself gives
\begin{equation}
 \xi(s)=\xi(1-s),\qquad
 \xi(\bar s)=\overline{\xi(s)},\qquad
 \xi(0)=\xi(1)=1/2.
 \label{eq:elem-xi-symmetries}
\end{equation}
These identities hold between the explicit approximants with their
error bounds, so an abstract identity theorem is unnecessary.

\subsection{Reciprocal gamma factors and their zeros}

\begin{lemma}[Constructive reciprocal gamma zero classification]
\label{lem:elem-gamma}
The entire reciprocal gamma function
\begin{equation}
 G(z)=z e^{\gamma z}\prod_{n=1}^\infty(1+z/n)e^{-z/n}
 \label{eq:elem-gamma-product}
\end{equation}
has precisely the simple zeros $0,-1,-2,\ldots$, and
$G'(-n)=(-1)^n n!$.  On a rational bounded region with supplied positive
separations from these integers, the product supplies an explicit
positive lower bound for $|G|$.
\end{lemma}

\begin{proof}
For $|z|\le R$ and an integer $N>2R$, the logarithm of the tail of
\eqref{eq:elem-gamma-product} satisfies
\begin{equation}
 \left|\sum_{n>N}\bigl(\log(1+z/n)-z/n\bigr)\right|
 \le2R^2\sum_{n>N}n^{-2}\le2R^2/N.
 \label{eq:elem-gamma-tail}
\end{equation}
This uses only the power series for the logarithm in $|z/n|\le1/2$;
there is no choice of a global logarithm branch.  The tail is the
exponential of a number of bounded modulus and is therefore apart from
zero.  Every finite initial factor whose linear term is separated from
zero is apart from zero as well.

The classification does not infer a disjunction from a product being
zero.  Here is a bounded procedure which avoids that invalid general
constructive inference.  Approximate $z$ by a rational complex number
$z_0$ with error at most $1/100$.  Search the finite list
$0,-1,\ldots,-N$, where $N>2R$, for a center $-n$ with
$|z_0+n|\le1/4$, using squared distances.  If there is no such center,
every relevant factor has modulus at least $1/5$, after its explicit
integer denominator is included.  The finite product and
\eqref{eq:elem-gamma-tail} then certify $G(z)$ apart from zero.  If
there is such a center, it is unique, and $|z+n|<1/3$.  All other
linear factors are separated from zero by at least $2/3$ before their
integer denominators are included.  Thus
$G(z)=(z+n)H_n(z)$ on this disk, where $H_n$ has an explicit positive
lower bound.  Consequently $G(z)=0$ implies $z=-n$.

For example, off all the disks, the lower bound is a finite product of
$1/5$, the numbers $1/(5j)$ for $1\le j\le N$, and a rational lower
bound for $\exp(-R(N+1)-2R^2/N)$.  Such a lower bound follows from
$e<4$ by an integer power of $1/4$.  The disk case omits the one
distinguished factor and is otherwise identical.

For completeness, these normalizations are obtained directly from the
finite products.  Put $H_N=\sum_{j=1}^N1/j$.  At $z=1$ the finite
product is $e^\gamma(N+1)e^{-H_N}$, which tends to one with the
explicit harmonic-sum estimate
$|H_N-\log N-\gamma|\le1/N$.  If $G_N$ denotes the finite product,
cancellation after multiplication, as a polynomial identity, gives
\[
 zG_N(z+1)=(N+z+1)e^{\gamma-H_N}G_N(z).
\]
The factor on the right tends to one with a uniform error on every
supplied bounded region.  Passing to the explicit limits gives the
entire identity $G(z)=zG(z+1)$ without dividing by a possibly zero
factor.
Thus $G(1)=1$ and the shift identity have explicit finite proofs.
Iterating the shift through the finitely many factors gives
\[
 G(z)=z(z+1)\cdots(z+n)G(z+n+1).
\]
At $z=-n$, the derivative is $(-1)^n n!$.  Its positive magnitude
proves simplicity.
\end{proof}

Away from $s=0,1$ the continuation is equivalently
\begin{equation}
 \zeta(s)=\frac{2\pi^{s/2}\xi(s)G(s/2)}{s(s-1)}.
 \label{eq:elem-zeta-from-xi}
\end{equation}
This identity agrees with the right-half-plane definition and extends
it by the effective constructions.  In particular it is not an
additional nonvanishing assumption about gamma.

\subsection{Trivial zeros and strip localization}

\begin{theorem}[The trivial zeros are simple]
\label{thm:elem-trivial}
In $\Re s<0$, the zeros of $\zeta$ are exactly $s=-2n$, $n\ge1$.
They are simple, with
\begin{equation}
 \zeta'(-2n)=(-1)^n
 \frac{(2n)!}{2(2\pi)^{2n}}\zeta(2n+1),\qquad
 |\zeta'(-2n)|\ge\frac{(2n)!}{2\,8^{2n}}>0.
 \label{eq:elem-trivial-derivative}
\end{equation}
The disk $|s+2n|\le1/4$ contains no other zero.
\end{theorem}

\begin{proof}
For $\Re s<0$, the point $1-s$ lies in $\Re s>1$.
Theorem~\ref{thm:elem-euler}, the gamma bounds, and
\eqref{eq:elem-xi-symmetries} show that $\xi(s)=\xi(1-s)$ is apart
from zero.  Every other factor in \eqref{eq:elem-zeta-from-xi} except
$G(s/2)$ is apart from zero.  Lemma~\ref{lem:elem-gamma} therefore
gives precisely the listed zeros, and its local factorization gives
their simplicity and isolation on the stated disks.

For the derivative, differentiate just the vanishing factor in
\eqref{eq:elem-zeta-from-xi}.  Since $G'(-n)=(-1)^n n!$,
\[
 \zeta'(-2n)
 =\frac{(-1)^n n!\,\pi^{-n}\xi(1+2n)}{2n(2n+1)}.
\]
Substitute the right-half-plane expression for $\xi(1+2n)$ and use
\[
 \Gamma(n+1/2)=\frac{(2n)!}{4^n n!}\sqrt\pi.
\]
This last identity is obtained by $n$ applications of the gamma
recurrence to $\Gamma(1/2)=\sqrt\pi$; a general gamma reflection
formula is unnecessary.  The resulting expression is
\eqref{eq:elem-trivial-derivative}.  Finally $2\pi<8$ and
$\zeta(2n+1)\ge1$ give its rational positive lower bound.
\end{proof}

At $s=0$, the cancellation $G(s/2)/s\longrightarrow1/2$ in
\eqref{eq:elem-zeta-from-xi} gives $\zeta(0)=-1/2$.
Theorem~\ref{thm:elem-F-clearance} and reflection show that $\xi$ is
apart from zero on the imaginary axis as well as on the one-line.
The reciprocal gamma factor has no additional zero there.  These
observations exclude the boundary from the nontrivial zero set.

\begin{theorem}[Quantitative critical-strip localization]
\label{thm:elem-strip}
Define a nontrivial zero as a zero of the entire completion $\xi$.
If $\xi(s)=0$ and a supplied integer $H\ge1$ satisfies $|\Im s|\le H$,
then, with $\varepsilon_H=2^{-40}(H+1)^{-5}$,
\begin{equation}
 \frac{\varepsilon_H}{2}<\Re s
 <1-\frac{\varepsilon_H}{2}.
 \label{eq:elem-quantitative-strip}
\end{equation}
Inside this strip, the zeros and multiplicities of $\xi$ and $\zeta$
coincide.
\end{theorem}

\begin{proof}
In the band $|\Re s-1|\le\varepsilon_H$, write
\[
 \xi(s)=\tfrac12s\pi^{-s/2}\Gamma(s/2)F(s).
\]
The factors are apart from zero by the gamma product bounds and
Theorem~\ref{thm:elem-F-clearance}.  In the remaining right-hand
region $\Re s\ge1+\varepsilon_H/2$, the Euler lower bound applies,
again with apartness of every other factor.  The overlapping regions
thus give apartness whenever $\Re s\ge1-\varepsilon_H$.

To obtain the strict conclusion rather than only its double negation,
approximate $\Re s$ within $\varepsilon_H/8$ and compare the rational
answer with $1-3\varepsilon_H/4$.  The larger alternative places $s$
in the certified nonzero region and contradicts $\xi(s)=0$.  The
smaller alternative supplies $\Re s<1-\varepsilon_H/2$ directly.
Reflection $s\mapsto1-s$ gives the lower bound.  Finally
\eqref{eq:elem-zeta-from-xi} has an analytic invertible factor on
$0<\Re s<1$, proving the multiplicity assertion.
\end{proof}

\subsection{The real axis and an explicit initial zero-free rectangle}

\begin{lemma}[Positive real-axis completion]
\label{lem:elem-real-positive}
For every real $x$, $\xi(x)\ge3/8$.  For a positive integer $R$ and
$|s|\le R$,
\begin{equation}
 |\xi'(s)|\le D_R
 :=2(2R+1)R!+R(R+1)(R+1)!.
 \label{eq:elem-xi-derivative-bound}
\end{equation}
In particular $D_2=56$ and
\begin{equation}
 |\xi(s)|\ge1/4\quad
 (0\le\Re s\le1,\quad |\Im s|\le1/448).
 \label{eq:elem-initial-rectangle}
\end{equation}
\end{lemma}

\begin{proof}
For $u\ge1$, $\pi>3$ and $n^2\ge n$ give
\[
 0\le\psi(u)\le\sum_{n\ge1}e^{-3nu}
 =\frac{e^{-3u}}{1-e^{-3u}}\le2e^{-3u}.
\]
For real $x$ outside $(0,1)$, the correction in
\eqref{eq:elem-xi-integral} is nonnegative, so $\xi(x)\ge1/2$.
For $0\le x\le1$, each of the two powers of $u$ in $J(x)$ is at
most one.  Hence
\[
 0\le J(x)\le4\int_1^\infty e^{-3u}\,du
 =\tfrac43e^{-3}<1.
\]
Since $x(x-1)\ge-1/4$, we obtain $\xi(x)\ge1/2-1/8=3/8$.
These inequalities can first be checked for rational $x$, where the
case distinction is decidable, and then passed with their uniform
continuity estimates to a named real.  No decision about the sign of
an arbitrary real is required.

For $|s|\le R$, both complex powers in $J(s)$ have modulus at most
$u^R$.  Using $e^{-3u}\le e^{-u}$ and $\log u\le u$ gives
\[
 |J(s)|\le4\int_0^\infty u^Re^{-u}\,du=4R!,
 \qquad |J'(s)|\le2(R+1)!.
\]
Differentiate \eqref{eq:elem-xi-integral} and use
$|2s-1|\le2R+1$ and $|s(s-1)|\le R(R+1)$ to obtain
\eqref{eq:elem-xi-derivative-bound}.  The factorial integrals are
proved by a finite number of integrations by parts with an explicit
tail.

The vertical segment from $\sigma$ to $\sigma+it$, for
$0\le\sigma\le1$ and $|t|\le1/448$, stays inside $|s|\le2$.
The derivative bound therefore yields
\[
 |\xi(\sigma+it)|\ge\xi(\sigma)-56|t|
 \ge3/8-56/448=1/4.
\]
This proves \eqref{eq:elem-initial-rectangle}.
\end{proof}

\begin{corollary}[Symmetry, with multiplicity]
\label{cor:elem-symmetry}
If $\rho$ is a nontrivial zero of multiplicity $m$, then
$\bar\rho$, $1-\rho$, and $1-\bar\rho$ are zeros of the same
multiplicity.  There are no nontrivial real zeros.
\end{corollary}

\begin{proof}
The two identities in \eqref{eq:elem-xi-symmetries} carry the local
Taylor expansion at $\rho$ to those at the three listed points and
preserve the first nonzero coefficient.  Equivalently they carry the
local factorization and its positive leading-coefficient margin to the
reflected factorizations.  The real-axis assertion is
Lemma~\ref{lem:elem-real-positive}.  Some of the listed points may
coincide; the conclusion does not assume four distinct zeros.
\end{proof}

\subsection{Arithmetic content of the estimates}

For each supplied integer height bound, all constants above are
rational expressions in integer powers and factorials.  The Euler,
prime-power, and gamma products have explicit primitive recursive
truncation lengths.  The horizontal and vertical integral bounds use
finite quadrature estimates, not a choice of an intermediate point.
The few distinctions involving named reals are implemented by
overlapping rational comparisons.  Thus the statements have finite
arithmetical versions with no hidden unbounded search.  Their full
object-theory formalization still requires the derivations of the
finite analytic identities and error bounds in the chosen iPRA
calculus; an ambient Lean theorem about classical complex numbers is
not a substitute for those derivations.
